\documentclass[11pt]{article}
\usepackage[utf8]{inputenc}
\usepackage[T1]{fontenc}
\usepackage{lmodern}
\usepackage[left=2.5cm, right=2.5cm, top=2.5cm, bottom=2.5cm]{geometry}
\usepackage{graphicx}
\usepackage{float}
\usepackage{xcolor}
\usepackage{booktabs}
\usepackage{tabularx}
\usepackage{longtable}
\usepackage{array}
\usepackage{hyperref}
\usepackage{listings}
\usepackage{caption}
\usepackage{amsmath}
\usepackage{enumitem}
\definecolor{MyPurple}{RGB}{120,0,180}
\definecolor{LightGray}{gray}{0.975}
\definecolor{CodeFrame}{RGB}{210,210,210}
\definecolor{SqlBlue}{RGB}{31,98,135}
\hypersetup{colorlinks=true, linkcolor=MyPurple, urlcolor=blue, citecolor=MyPurple}
\setlength{\parindent}{0pt}
\setlength{\parskip}{5pt}
\renewcommand{\arraystretch}{1.18}
\lstdefinestyle{sqlstyle}{
    language=SQL,
    basicstyle=\ttfamily\scriptsize,
    keywordstyle=\color{SqlBlue}\bfseries,
    commentstyle=\color{gray},
    stringstyle=\color{MyPurple},
    breaklines=true,
    breakatwhitespace=false,
    showstringspaces=false,
    frame=single,
    rulecolor=\color{CodeFrame},
    backgroundcolor=\color{LightGray},
    columns=fullflexible,
    keepspaces=true,
    tabsize=2
}
\lstdefinestyle{makestyle}{
    basicstyle=\ttfamily\scriptsize,
    breaklines=true,
    frame=single,
    rulecolor=\color{CodeFrame},
    backgroundcolor=\color{LightGray},
    columns=fullflexible,
    keepspaces=true
}
\newcommand{\dictheader}{\hline \textbf{Attribute} & \textbf{Type} & \textbf{Constraints} & \textbf{Description} \\ \hline}
\title{Databases\\\textbf{Final Report: Designing a Telemetry and UX Database for \textit{Chocolate-Doom}}}
\author{Fergie Lindsay D\'iaz, Juan Esteban Ortega, Juan Esteban, Federico Mejia}
\date{June 2026}
\begin{document}
\maketitle
\tableofcontents
\clearpage

\section{Introduction and context} 

This project is about designing and building a database to store and analyze data collected from a modified version of the game \textit{\textbf{Chocolate-Doom}}. The game records detailed information during gameplay, such as the player’s position, movement, and status at every moment. In addition to this, players also provide basic personal information and complete a user experience survey. 

\medskip The main goal in the present report or document of this project is to organize all this data in a structured way so it can be stored, accessed, and analyzed efficiently. To do this, we will design a relational database that connects gameplay data with user and survey information. The database should be able to handle a large amount of data and ensure that everything is consistent and well organized. 

\medskip Another important part of the project is to run queries that help us understand player behavior. For example, we will look at how each player moves, which sectors are more visited,
and how individual trajectories develop during the game. We will also relate this information to their survey responses to see how their experience connects to their behavior. 

\medskip Finally, we will follow good practices for data handling, including cleaning and loading the data properly, and considering basic privacy and ethical aspects. Overall, the project aims to combine database design with practical analysis of real gameplay data.



\section{Assumptions and requirements}

\subsection{Assumptions}
\begin{itemize}
  \item Each User is a real participant who has accepted consent before any data is collected.
  \item A User can have one or more Player identities (nicknames), but each Player belongs to only one User.
  \item Each Game represents one gameplay session with a start and end time.
  \item Each Game is linked to exactly one Player.
  \item Telemetry data is recorded at every tic (very frequent events) for each game session.
  \item Each telemetry record is uniquely identified by (game\_id, tic).
  \item Each Game is linked to exactly one Map and one Episode.
  \item Maps are divided into Sectors, which help us analyze player positions.
  \item Each user is expected to submit one UX survey after gameplay.
  \item Survey responses are completed after the game session and linked to the user.
\end{itemize}


\subsection{Functional requirements}

\begin{itemize}
  \item Store and manage user demographic information (age, gender, experience level, consent).
  \item Support multiple player identities per user.
  \item Record high-frequency gameplay telemetry including position, momentum, orientation, and combat stats.
  \item Store game session metadata including player, episode, map, start time, and end time.
  \item Support spatial grouping of telemetry data by map sectors.
  \item Ingest raw telemetry logs (TSV format) into a structured relational database.
  \item Store and manage UX survey instruments (PENS, GUESS, BANGS).
  \item Store user responses to UX instruments at the question level.
  \item Support analytical queries on movement, trajectories, and player behavior.
  \item Support aggregation of telemetry data for trajectory and gameplay behavior analysis.
\end{itemize}

\subsection{Non-functional requirements}
\begin{itemize}
  \item The database should handle large telemetry datasets efficiently (at least 20,000 records, but ideally much more as data grows).
  \item The schema should work well as the number of players, games, and telemetry records increases, without needing major changes.
  \item Primary keys and foreign keys should be used to keep data consistent across all tables.
  \item Indexes should be added to improve performance for common queries, especially those involving time sequences and spatial data.
  \item The design should allow new telemetry fields or new UX instruments to be added without redesigning the whole database.
  \item Data loading should be repeatable using a clear ETL process so the database can be rebuilt from raw logs if needed.
  \item User data should be anonymized, and consent must be stored before any data is collected or used.
  \item If some telemetry records are invalid or broken, they should be logged and skipped without stopping the whole ingestion process.
\end{itemize}

\clearpage \section{Survey Election - GUESS}

The survey we decided to work with and base our design on is Game User-Experience Satisfaction Scale \href{https://uxpajournal.org/wp-content/uploads/sites/7/pdf/JUS_Keebler_GUESS-18%20Scoring%20Guidelines.pdf}{\textbf{GUESS}}. 

\subsection{What is it?}

\medskip  The Game User Experience Satisfaction Scale (GUESS-18) is a survey used to understand how players feel about a video game. It has 18 questions that are divided into nine sections, and each one looks at a different part of the experience, like how easy the game is to use, how fun it is, how immersive it feels, or how good the visuals and audio are.

\medskip  Each question is answered using a scale from 1 to 7, where players say how much they agree or disagree. At the end, the answers are grouped to get scores for each section and also a total score that shows the overall satisfaction. There is one question that is reverse scored (“I feel bored while playing the game”), so it needs to be handled a bit differently.

\medskip  GUESS-18 is actually a shorter version of a longer survey with 55 questions, but this one is faster to complete and still gives good results.

\subsection{Why did we chose it?}

\medskip We chose GUESS-18 because it is easy to use and still gives a lot of useful information. It lets us measure different aspects of the player experience, not just if the game works, but also if it is fun, engaging, or interesting.

\medskip Another reason is that it is a well-known and tested survey, so we can trust the results we get. This is important because we want to compare what players say with the data we collect from the game.

\medskip It also fits nicely with our database design. Since the survey is divided into sections and questions, we can store everything in an organized way and later analyze things like enjoyment or immersion separately.

\medskip Finally, it is short, so players can answer it quickly after playing. This makes it more likely that we get complete and honest responses.

\clearpage

\section{Entity-Relation Diagram}

The diagram organizes gameplay telemetry, user information, and survey responses into coherent relationships. It highlights how users are linked to players, how players are linked to individual game sessions, and how each session generates detailed telemetry data. It also integrates UX instruments to connect behavioral data with subjective user experience.


\begin{figure}[H]
    \centering
    \includegraphics[width=1\linewidth]{imagen.png}
    \caption{Entity-Relation Diagram}
    \label{fig:placeholder}
\end{figure}

The diagram organizes gameplay telemetry, user information, and survey responses into coherent relationships. It highlights how users are linked to players, how players participate in game sessions, and how each session generates detailed telemetry data. It also integrates UX instruments to connect behavioral data with subjective user experience.



\section{Relational Schema}

The database is normalized by separating participant identity, gameplay metadata, map structure, telemetry events, and UX responses. High-frequency telemetry is isolated in \texttt{TelemetryEvent}, while relatively stable reference data such as \texttt{Episode}, \texttt{Map}, \texttt{Sector}, and UX metadata remain in separate tables.

The most important correction in the final schema is the composite key in \texttt{Game}. A simple primary key on \texttt{game\_id} would fail with the real telemetry because one session can include more than one participant. The final model therefore uses:

\begin{lstlisting}[style=sqlstyle]
PRIMARY KEY (game_id, player_id)
\end{lstlisting}

and \texttt{TelemetryEvent} includes \texttt{player\_id} and references the composite key:

\begin{lstlisting}[style=sqlstyle]
FOREIGN KEY (game_id, player_id)
REFERENCES Game(game_id, player_id)
ON DELETE CASCADE
\end{lstlisting}


\subsection{Core entities}
This subsection lists the core entities identified in the system. These represent the main objects stored in the database, including users, gameplay sessions, telemetry records, and UX survey components.

\begin{itemize}
    \item \texttt{AppUser}: participant demographics and informed consent.
    \item \texttt{Player}: in-game alias linked to a user.
    \item \texttt{Episode}, \texttt{Map}, \texttt{Sector}: level hierarchy and spatial grouping.
    \item \texttt{Game}: player participation in a shared session.
    \item \texttt{TelemetryEvent}: atomic per-tic player state.
    \item \texttt{UXInstrument}, \texttt{UXSubscale}, \texttt{UXItem}: UX questionnaire metadata.
    \item \texttt{UXResponse}, \texttt{UXResponseItem}: user responses at response and item level.
\end{itemize}

\subsection{Foreign Keys}

\smallskip
The way in which we write or show the foreign keys will later help us in the way we can use it in SQL.

\smallskip
\textit{\textbf{Example}}:
\begin{lstlisting}[style=sqlstyle]
          FOREIGN KEY (user_id) REFERENCES User(user_id)
\end{lstlisting}


This subsection defines the relationships between tables using foreign keys. These constraints enforce referential integrity and ensure consistency across users, gameplay sessions, telemetry events, and UX responses.


\clearpage
\subsection{Data Dictionary}

\subsubsection*{AppUser}
\begin{tabularx}{\linewidth}{|l|l|l|X|}
\hline
\textbf{Attribute} & \textbf{Type} & \textbf{Constraints} & \textbf{Description} \\
\hline
user\_id & SERIAL & PK & Unique identifier for each user \\
\hline
age & INTEGER & CHECK (10--100) & Age of the participant \\
\hline
gender & VARCHAR(50) & --- & Self-reported gender \\
\hline
experience\_level & VARCHAR(50) & --- & Level of gaming experience \\
\hline
consent & BOOLEAN & NOT NULL, DEFAULT FALSE & Indicates informed consent for participation \\
\hline
\end{tabularx}

\subsubsection*{Player}
\begin{tabularx}{\linewidth}{|l|l|l|X|}
\hline
player\_id & SERIAL & PK & Unique identifier for each in-game player \\
\hline
user\_id & INTEGER & FK, NOT NULL & References AppUser(user\_id); links player to a user \\
\hline
nickname & VARCHAR(100) & NOT NULL & Player's in-game alias \\
\hline
\end{tabularx}

\subsubsection*{Episode}
\begin{tabularx}{\linewidth}{|l|l|l|X|}
\hline
episode\_id & SERIAL & PK & Unique episode identifier \\
\hline
name & VARCHAR(100) & NOT NULL & Name of the episode \\
\hline
\end{tabularx}

\subsubsection*{Map}
\begin{tabularx}{\linewidth}{|l|l|l|X|}
\hline
map\_id & SERIAL & PK & Unique map identifier \\
\hline
episode\_id & INTEGER & FK & References Episode(episode\_id) \\
\hline
map\_code & VARCHAR(50) & NOT NULL & Map code (e.g., E1M1) \\
\hline
\end{tabularx}

\subsubsection*{Sector}
\begin{tabularx}{\linewidth}{|l|l|l|X|}
\hline
sector\_id & SERIAL & PK & Unique sector identifier \\
\hline
map\_id & INTEGER & FK & References Map(map\_id) \\
\hline
grid\_x & INTEGER & --- & X-coordinate within the map grid \\
\hline
grid\_y & INTEGER & --- & Y-coordinate within the map grid \\
\hline
\end{tabularx}

\subsubsection*{Game}
\begin{tabularx}{\linewidth}{|l|l|l|X|}
\hline
game\_id & SERIAL & PK & Unique identifier for each game session \\
\hline
player\_id & INTEGER & FK & References Player(player id)\\
\hline
map\_id & INTEGER & FK & References Map(map\_id) \\
\hline
start\_time & TIMESTAMP & NOT NULL & Timestamp when the session started \\
\hline
end\_time & TIMESTAMP & --- & Timestamp when the session ended \\
\hline
\end{tabularx}

\subsubsection*{TelemetryEvent}
\begin{tabularx}{\linewidth}{|l|l|l|X|}
\hline
event\_id & BIGSERIAL & PK & Unique identifier for each telemetry record \\
\hline
game\_id & INTEGER & FK & References Game(game\_id) \\
\hline
player\_id & INTEGER & FK & References Player(player\_id) \\
\hline
sector\_id & INTEGER & FK & Nullable reference to Sector(sector\_id) \\
\hline
tic & INTEGER & NOT NULL & Discrete time unit (game tick) \\
\hline
pos\_x & DOUBLE PRECISION & NOT NULL & Player position on X-axis \\
\hline
pos\_y & DOUBLE PRECISION & NOT NULL & Player position on Y-axis \\
\hline
mom\_x & DOUBLE PRECISION & --- & Momentum on X-axis \\
\hline
mom\_y & DOUBLE PRECISION & --- & Momentum on Y-axis \\
\hline
angle & DOUBLE PRECISION & --- & Player orientation angle \\
\hline
fov & DOUBLE PRECISION & --- & Field of view \\
\hline
health & INTEGER & --- & Player health value \\
\hline
armor & INTEGER & --- & Player armor value \\
\hline
ammo & INTEGER & --- & Ammunition count \\
\hline
\end{tabularx}

\textbf{Constraint:} UNIQUE (game\_id, tic)

\subsubsection*{UXInstrument}
\begin{tabularx}{\linewidth}{|l|l|l|X|}
\hline
instrument\_id & SERIAL & PK & Unique identifier for each UX instrument \\
\hline
name & VARCHAR(100) & NOT NULL & Name of the instrument (e.g., GUESS, PENS) \\
\hline
version & VARCHAR(50) & --- & Instrument version \\
\hline
scale\_min & INTEGER & --- & Minimum possible score \\
\hline
scale\_max & INTEGER & --- & Maximum possible score \\
\hline
\end{tabularx}

\subsubsection*{UXSubscale}
\begin{tabularx}{\linewidth}{|l|l|l|X|}
\hline
subscale\_id & SERIAL & PK & Unique identifier for each subscale \\
\hline
instrument\_id & INTEGER & FK & References UXInstrument(instrument\_id) \\
\hline
name & VARCHAR(100) & NOT NULL & Name of the subscale \\
\hline
code & VARCHAR(20) & --- & Short code identifier \\
\hline
\end{tabularx}

\subsubsection*{UXItem}
\begin{tabularx}{\linewidth}{|l|l|l|X|}
\hline
item\_id & SERIAL & PK & Unique identifier for each item/question \\
\hline
instrument\_id & INTEGER & FK & References UXInstrument(instrument\_id) \\
\hline
subscale\_id & INTEGER & FK & References UXSubscale(subscale\_id) \\
\hline
item\_number & INTEGER & NOT NULL & Item order within the instrument \\
\hline
question\_text & TEXT & NOT NULL & Text of the survey question \\
\hline
is\_reverse\_scored & BOOLEAN & DEFAULT FALSE & Indicates if reverse scoring is applied \\
\hline
\end{tabularx}

\subsubsection*{UXResponse}
\begin{tabularx}{\linewidth}{|l|l|l|X|}
\hline
response\_id & SERIAL & PK & Unique identifier for each survey response \\
\hline
user\_id & INTEGER & FK & References AppUser(user\_id) \\
\hline
instrument\_id & INTEGER & FK & References UXInstrument(instrument\_id) \\
\hline
responded\_at & TIMESTAMP & DEFAULT CURRENT\_TIMESTAMP & Date and time of submission \\
\hline
\end{tabularx}

\subsubsection*{UXResponseItem}
\begin{tabularx}{\linewidth}{|l|l|l|X|}
\hline
response\_id & INTEGER & PK, FK & References UXResponse(response\_id) \\
\hline
item\_id & INTEGER & PK, FK & References UXItem(item\_id) \\
\hline
score & INTEGER & CHECK ($\geq 0$) & Score given to the item \\
\hline
\end{tabularx}



\clearpage
\section{Data Definition Language (DDL)}

\subsection{SQL Dictionary}

\begin{itemize}
  \item \href{https://www.geeksforgeeks.org/postgresql/postgresql-serial/}{\textbf{SERIAL - BIGSERIAL}}: Auto-increment numbers for IDs.  
    \textit{Example:} In \texttt{AppUser}, \texttt{user\_id SERIAL PRIMARY KEY} gives each user a unique ID. In \texttt{TelemetryEvent}, \texttt{event\_id BIGSERIAL PRIMARY KEY} allows billions of events.

  \item \href{https://www.w3schools.com/sql/sql_check.asp}{\textbf{CHECK}}: Rule that makes sure values are valid.  
    \textit{Example:} In \texttt{AppUser}, \texttt{age INTEGER CHECK (age BETWEEN 10 AND 100)} only allows ages between 10 and 100. In \texttt{UXResponseItem}, \texttt{score INTEGER CHECK (score >= 0)} ensures scores are not negative.

  \item \href{https://www.geeksforgeeks.org/sql/cascade-in-sql/}{\textbf{ON DELETE CASCADE}}: If a parent row is deleted, linked child rows are deleted too.  
    \textit{Example:} If an \textit{AppUser} is deleted, their \textit{Player} record is also deleted. If a Game is deleted, all its \textit{TelemetryEvent} records are removed.

  \item \href{https://www.w3resource.com/PostgreSQL/date-time-functions-and-operators.php}{\textbf{TIMESTAMP}}: Stores both date and time.  
    \textit{Example:} In \texttt{Game}, \texttt{start\_time TIMESTAMP NOT NULL} records when the game started. In \texttt{UXResponse}, \texttt{responded\_at TIMESTAMP DEFAULT CURRENT\_TIMESTAMP} saves the time a user answered.

  \item \href{https://www.postgresql.org/docs/current/datatype-numeric.html}{\textbf{PRECISION}}: Controls how exact numbers are stored.  
    \textit{Example:} In \texttt{TelemetryEvent}, \texttt{pos\_x DOUBLE PRECISION} stores very detailed player positions, momentum, and angles.
    
  \item \href{https://www.postgresql.org/docs/current/ddl-constraints.html}{\textbf{ON DELETE SET NULL}}: If a parent row is deleted, the link is removed but the child row stays (foreign key becomes \texttt{NULL}).  
    \textit{Example:} If a \texttt{Sector} is deleted, the \texttt{TelemetryEvent} still exists but its \texttt{sector\_id} becomes \texttt{NULL}. If a \texttt{UXSubscale} is deleted, the \texttt{UXItem} remains but loses its subscale link.
\end{itemize}

\subsection{Data Definition Language highlights}
The DDL was implemented in PostgreSQL. The final schema uses standard relational constraints and keeps telemetry fields typed as numeric values so that spatial and temporal analysis can be performed directly in SQL.

\begin{lstlisting}[style=sqlstyle]
CREATE TABLE Game (
    game_id INTEGER NOT NULL,
    player_id INTEGER NOT NULL REFERENCES Player(player_id) ON DELETE CASCADE,
    map_id INTEGER REFERENCES Map(map_id) ON DELETE CASCADE,
    start_time TIMESTAMP NOT NULL,
    end_time TIMESTAMP,
    PRIMARY KEY (game_id, player_id)
);

CREATE TABLE TelemetryEvent (
    event_id BIGSERIAL PRIMARY KEY,
    game_id INTEGER NOT NULL,
    player_id INTEGER NOT NULL,
    sector_id INTEGER REFERENCES Sector(sector_id) ON DELETE SET NULL,
    tic INTEGER NOT NULL,
    pos_x DOUBLE PRECISION NOT NULL,
    pos_y DOUBLE PRECISION NOT NULL,
    mom_x DOUBLE PRECISION,
    mom_y DOUBLE PRECISION,
    angle DOUBLE PRECISION,
    fov DOUBLE PRECISION,
    health INTEGER,
    armor INTEGER,
    ammo INTEGER,
    FOREIGN KEY (game_id, player_id)
        REFERENCES Game(game_id, player_id)
        ON DELETE CASCADE,
    CONSTRAINT uq_telem UNIQUE (game_id, player_id, tic)
);
\end{lstlisting}

\section{SQL CODES}

\subsection{parte2\_puntoB.sql}

In this section, we explain how the telemetry data is loaded into the database.

\smallskip
First, the raw data from the TSV file is stored in the \textit{staging\_telemetry} table. In this table, all values are saved as text to make the loading process simple.

\smallskip
Then, we check for invalid records (for example, missing game, player, or tic information) and store them in the \textit{ingestion\_error\_log }table.

\smallskip

Finally, the valid data is cleaned, converted to the correct data types, and inserted into the \texttt{TelemetryEvent} table. Duplicate records are avoided, and empty values are handled properly.

\begin{lstlisting}[style=sqlstyle]
-- 1. Crear tabla de Staging
CREATE TABLE IF NOT EXISTS staging_telemetry (
    game_id_raw      TEXT,
    player_id_raw    TEXT,
    sector_id_raw    TEXT,
    tic_raw          TEXT,
    pos_x_raw        TEXT,
    pos_y_raw        TEXT,
    mom_x_raw        TEXT,
    mom_y_raw        TEXT,
    angle_raw        TEXT,
    fov_raw          TEXT,
    health_raw       TEXT,
    armor_raw        TEXT,
    ammo_raw         TEXT
);

-- 2. Crear tabla de Log de Errores 
CREATE TABLE IF NOT EXISTS ingestion_error_log (
    error_id   SERIAL PRIMARY KEY,
    game_id    TEXT,
    tic        TEXT,
    motivo     TEXT,
    fecha_err  TIMESTAMP DEFAULT CURRENT_TIMESTAMP
);

-- 3. Proceso ETL 
INSERT INTO ingestion_error_log (game_id, tic, motivo)
SELECT game_id_raw, tic_raw, 
FROM staging_telemetry
WHERE game_id_raw IS NULL OR player_id_raw IS NULL OR tic_raw IS NULL;

INSERT INTO TelemetryEvent (
    game_id, player_id, sector_id, tic, 
    pos_x, pos_y, 
    mom_x, mom_y, 
    angle, fov, health, armor, ammo
)
SELECT DISTINCT ON (CAST(game_id_raw AS INT), CAST(player_id_raw AS INT), CAST(tic_raw AS INT))
    CAST(game_id_raw AS INT),
    CAST(player_id_raw AS INT),
    CAST(NULLIF(sector_id_raw, '') AS INT),
    CAST(tic_raw AS INT),
    CAST(pos_x_raw AS DOUBLE PRECISION),
    CAST(pos_y_raw AS DOUBLE PRECISION),
    CAST(NULLIF(mom_x_raw, '') AS DOUBLE PRECISION),
    CAST(NULLIF(mom_y_raw, '') AS DOUBLE PRECISION),
    CAST(NULLIF(angle_raw, '') AS DOUBLE PRECISION),
    CAST(NULLIF(fov_raw, '') AS DOUBLE PRECISION),
    CAST(NULLIF(health_raw, '') AS INT),
    CAST(NULLIF(armor_raw, '') AS INT),
    CAST(NULLIF(ammo_raw, '') AS INT)
FROM staging_telemetry
WHERE game_id_raw IS NOT NULL 
  AND player_id_raw IS NOT NULL 
  AND tic_raw IS NOT NULL
ORDER BY CAST(game_id_raw AS INT), CAST(player_id_raw AS INT), CAST(tic_raw AS INT);
\end{lstlisting}

\clearpage
\subsection{parte\_3\_puntoB.sql}

In this section we insert the GUESS-18 survey structure into the database, including the instrument, its subscales, and all items.

\medskip

The field \texttt{is\_reverse\_scored} is used to indicate whether a question needs reverse scoring:

\medskip
\begin{itemize}
    \item \textbf{TRUE:} The item is reverse scored. This means that higher values represent a negative response, so the score must be inverted during analysis.

\medskip

    \item \textbf{FALSE:} The item is not reverse scored. The values can be interpreted directly, where higher scores indicate a more positive experience.

\end{itemize}


The questionnaire is divided into subscales, which are groups of questions that measure different aspects of the player experience.

Each subscale focuses on a specific category, such as :
\begin{itemize}
    \item \textbf{Usability/Playability} → how easy the game is to play
    \item \textbf{Narratives} → story and engagement with the game's plot
    \item \textbf{Play Engrossment} → level of immersion and focus during gameplay
    \item \textbf{Enjoyment} → how fun the game feels
    \item \textbf{Creative Freedom} → ability to be imaginative while playing
    \item \textbf{Audio Aesthetics} → quality and impact of sound and music
    \item \textbf{Personal Gratification} → focus on performance and achievement
    \item \textbf{Visual Aesthetics} → visual quality and appeal of the game
\end{itemize}

This organization allows the analysis to evaluate different dimensions of the game experience separately, instead of treating all responses as a single score.

In our implementation, the GUESS-18 survey is divided into nine subscales, each containing two items or questions, but we used \textbf{8 / 9} because one of these included the experience with other players, what was discussed and taken out to follow instructions and rules of the game: one player.

Each question (item) belongs to one subscale, allowing us to group responses and analyze different aspects of the player experience independently.

Each subscale is stored in the \texttt{UXSubscale} table and is linked to its corresponding questions through the \texttt{UXItem} table. This structure allows efficient organization and analysis of survey data.

\begin{itemize}
    \item Subscale \textbf{=} Topic of analysis for the survey.
    \item Item \textbf{=} Question
\end{itemize}

\clearpage
\begin{lstlisting}[style=sqlstyle]
-- 1. Insertar el instrumento GUESS-18
INSERT INTO UXInstrument (instrument_id, name, version, scale_min, scale_max) 
VALUES (1, 'GUESS', '18-Item Short Scale', 1, 7);
 
-- 2. Insertar 8/9 Subescalas oficiales del GUESS-18
INSERT INTO UXSubscale (subscale_id, instrument_id, name, code) VALUES 
(1, 1, 'Usability/Playability', 'USAB'),
(2, 1, 'Narratives', 'NARR'),
(3, 1, 'Play Engrossment', 'ENGR'),
(4, 1, 'Enjoyment', 'ENJO'),
(5, 1, 'Creative Freedom', 'CREA'),
(6, 1, 'Audio Aesthetics', 'AUDI'),
(7, 1, 'Personal Gratification', 'GRAT'),
(8, 1, 'Visual Aesthetics', 'VISU');
 
-- 3. Insertar los 16 ítems (2 por subescala)
INSERT INTO UXItem (item_id, instrument_id, subscale_id, item_number, question_text, is_reverse_scored) VALUES 
-- Usability/Playability
(1, 1, 1, 1, 'I find the controls of the game to be straightforward.', FALSE),
(2, 1, 1, 2, 'I find the game''s interface to be easy to navigate.', FALSE),
 
-- Narratives
(3, 1, 2, 3, 'I am captivated by the game''s story from the beginning.', FALSE),
(4, 1, 2, 4, 'I enjoy the fantasy or story provided by the game.', FALSE),
 
-- Play Engrossment
(5, 1, 3, 5, 'I feel detached from the outside world while playing the game.', FALSE),
(6, 1, 3, 6, 'I do not care to check events that are happening in the real world during the game.', FALSE),
 
-- Enjoyment
(7, 1, 4, 7, 'I think the game is fun.', FALSE),
(8, 1, 4, 8, 'I feel bored while playing the game.', TRUE),
 
-- Creative Freedom
(9, 1, 5, 9, 'I feel the game allows me to be imaginative.', FALSE),
(10, 1, 5, 10, 'I feel creative while playing the game.', FALSE),
 
-- Audio Aesthetics
(11, 1, 6, 11, 'I enjoy the sound effects in the game.', FALSE),
(12, 1, 6, 12, 'I feel the game''s audio (e.g., sound effects, music) enhances my gaming experience.', FALSE),
 
-- Personal Gratification
(13, 1, 7, 13, 'I am very focused on my own performance while playing the game.', FALSE),
(14, 1, 7, 14, 'I want to do as well as possible during the game.', FALSE),
 
-- Visual Aesthetics
(17, 1, 9, 17, 'I enjoy the game''s graphics.', FALSE),
(18, 1, 9, 18, 'I think the game is visually appealing.', FALSE);
 
-- Resincronizar las secuencias en PostgreSQL para futuros INSERTS
SELECT setval('uxinstrument_instrument_id_seq', (SELECT MAX(instrument_id) FROM UXInstrument));
SELECT setval('uxsubscale_subscale_id_seq', (SELECT MAX(subscale_id) FROM UXSubscale));
SELECT setval('uxitem_item_id_seq', (SELECT MAX(item_id) FROM UXItem));
\end{lstlisting}

\clearpage
\section{C.3 Views and Materialized View}
This section defines SQL views and a materialized view used to support frequent analytical queries over gameplay and telemetry data. These structures simplify complex queries and improve performance for repeated analyses.

\subsection{View 1: Player Game Summary}

\begin{lstlisting}[style=sqlstyle]
CREATE VIEW player_game_summary AS
SELECT
    p.player_id,
    p.nickname,
    COUNT(g.game_id) AS total_games,
    MIN(g.start_time) AS first_game,
    MAX(g.end_time) AS last_game
FROM Player p
JOIN Game g ON p.player_id = g.player_id
GROUP BY p.player_id, p.nickname;
\end{lstlisting}

This view summarizes the activity of each player by counting the number of games played and showing their first and last recorded sessions.

\bigskip

\subsection{View 2: Map and Episode Activity Summary}

\begin{lstlisting}[style=sqlstyle]
CREATE VIEW map_activity_summary AS
SELECT
    e.episode_id,
    e.name AS episode_name,
    m.map_id,
    m.map_code,
    COUNT(g.game_id) AS total_games
FROM Episode e
JOIN Map m ON e.episode_id = m.episode_id
JOIN Game g ON m.map_id = g.map_id
GROUP BY e.episode_id, e.name, m.map_id, m.map_code;
\end{lstlisting}

This view provides a summary of gameplay activity per map and episode, allowing analysis of which levels are most frequently played.

\bigskip

\subsection{Materialized View: Player Telemetry Summary}

\begin{lstlisting}[style=sqlstyle]
CREATE MATERIALIZED VIEW player_telemetry_summary AS
SELECT
    g.game_id,
    g.player_id,
    COUNT(t.event_id) AS total_events,
    AVG(t.pos_x) AS avg_pos_x,
    AVG(t.pos_y) AS avg_pos_y,
    AVG(t.health) FILTER (WHERE t.health IS NOT NULL) AS avg_health,
    AVG(t.ammo) FILTER (WHERE t.ammo IS NOT NULL) AS avg_ammo
FROM Game g
JOIN TelemetryEvent t ON g.game_id = t.game_id
GROUP BY g.game_id, g.player_id;
\end{lstlisting}

This materialized view precomputes aggregated telemetry statistics per game session. Since telemetry data is large and frequently queried, storing pre-aggregated results improves performance significantly.

A materialized view is used instead of a regular view because regular views execute the query every time they are called. In contrast, the materialized view saves the results, making access faster. However, this comes with the limitation that the data must be manually refreshed whenever new telemetry data is added.

\subsection{Refreshing the Materialized View}

\begin{lstlisting}[style=sqlstyle]
REFRESH MATERIALIZED VIEW player_telemetry_summary;
\end{lstlisting}

This command is used to update the materialized view when new telemetry data is inserted into the database.
\clearpage
\section{Analytical queries and results}
This section presents six analytical queries selected from the project brief. They were adapted to the corrected schema, where a shared session is represented by \texttt{game\_id} and each player participation is represented by \texttt{(game\_id, player\_id)}.

\subsection{Query 1: average session duration per map}
This query computes the average duration of game participations per map.

\begin{lstlisting}[style=sqlstyle]
SELECT
    m.map_id,
    m.map_code,
    e.name AS episodio,
    COUNT(*) AS sesiones,
    AVG(g.end_time - g.start_time) AS duracion_promedio
FROM Game g
JOIN Map m     ON g.map_id = m.map_id
JOIN Episode e ON m.episode_id = e.episode_id
WHERE g.end_time IS NOT NULL
GROUP BY m.map_id, m.map_code, e.name
ORDER BY duracion_promedio DESC;
\end{lstlisting}

\begin{table}[H]
\centering
\begin{tabular}{llrr}
\toprule
\textbf{map code} & \textbf{episode} & \textbf{sessions} & \textbf{average duration} \\
\midrule
E1M1 & Knee-Deep in the Dead & 10 & 00:00:42.1 \\
E2M1 & The Shores of Hell & 10 & 00:00:31.3 \\
E3M1 & Inferno & 8 & 00:00:29.5 \\
\bottomrule
\end{tabular}
\caption{Average session duration by map.}
\end{table}

\textbf{Interpretation.} This query identifies which maps tend to generate longer gameplay sessions. Longer average durations may indicate higher engagement, greater map complexity, or larger exploration requirements.

\clearpage
\subsection{Query 2: players with the highest average proximity}
A self-join compares players at the same \texttt{game\_id} and \texttt{tic}. A lower average Euclidean distance indicates greater proximity.

\begin{lstlisting}[style=sqlstyle]
SELECT
    a.player_id AS jugador_a,
    b.player_id AS jugador_b,
    COUNT(*) AS tics_comparados,
    ROUND(AVG(sqrt(power(a.pos_x - b.pos_x, 2)
                 + power(a.pos_y - b.pos_y, 2)))::numeric, 2) AS distancia_promedio
FROM TelemetryEvent a
JOIN TelemetryEvent b
  ON a.game_id = b.game_id
 AND a.tic = b.tic
 AND a.player_id < b.player_id
GROUP BY a.player_id, b.player_id
ORDER BY distancia_promedio ASC;
\end{lstlisting}

\begin{table}[H]
\centering
\begin{tabular}{rrrr}
\toprule
\textbf{player A} & \textbf{player B} & \textbf{compared tics} & \textbf{avg. distance} \\
\midrule
3 & 4 & 1300 & 522.70 \\
4 & 6 & 1200 & 904.71 \\
2 & 4 & 1000 & 992.56 \\
2 & 6 & 1800 & 1093.84 \\
3 & 6 & 800 & 1165.15 \\
\bottomrule
\end{tabular}
\caption{Closest player pairs by average distance.}
\end{table}

\textbf{Interpretation.} Player pairs with lower average Euclidean distance tend to remain physically closer during gameplay. This may indicate cooperation, coordinated movement, or shared objectives.

\clearpage
\subsection{Query 3: shortest and longest trajectory distances per player}
This query uses \texttt{LAG()} to calculate step-by-step distance between consecutive tics, then aggregates trajectory length by player. Using a window function provides a clearer solution than repeatedly self-joining the table.

\begin{lstlisting}[style=sqlstyle]
WITH pasos AS (
    SELECT
        game_id, player_id, tic, pos_x, pos_y,
        LAG(pos_x) OVER w AS px_prev,
        LAG(pos_y) OVER w AS py_prev
    FROM TelemetryEvent
    WINDOW w AS (PARTITION BY game_id, player_id ORDER BY tic)
),
dist_por_partida AS (
    SELECT
        game_id,
        player_id,
        SUM(sqrt(power(pos_x - px_prev, 2) + power(pos_y - py_prev, 2))) AS dist_total
    FROM pasos
    WHERE px_prev IS NOT NULL
    GROUP BY game_id, player_id
)
SELECT
    player_id,
    COUNT(*) AS partidas_jugadas,
    ROUND(MIN(dist_total)::numeric, 2) AS trayectoria_min,
    ROUND(MAX(dist_total)::numeric, 2) AS trayectoria_max
FROM dist_por_partida
GROUP BY player_id
ORDER BY player_id;
\end{lstlisting}

\begin{table}[H]
\centering
\begin{tabular}{rrrr}
\toprule
\textbf{player} & \textbf{games} & \textbf{min distance} & \textbf{max distance} \\
\midrule
1 & 5 & 7587.75 & 17487.44 \\
2 & 5 & 8448.40 & 15889.34 \\
3 & 5 & 6060.07 & 11384.56 \\
4 & 5 & 5843.17 & 8502.11 \\
5 & 4 & 9838.55 & 18994.42 \\
6 & 4 & 6488.90 & 14039.98 \\
\bottomrule
\end{tabular}
\caption{Minimum and maximum trajectory distance by player.}
\end{table}

\textbf{Interpretation.} Trajectory length provides a quantitative measure of player movement. Large trajectory values suggest exploration-oriented behavior, while shorter trajectories may indicate defensive or objective-focused play styles.

\clearpage
\subsection{Query 4: UX responses for players with above-average trajectory duration}
Trajectory duration is measured as the number of telemetry tics per player. The global average is 5,750 tics. Players 1, 2, and 3 are above that average, so their UX scores are grouped by subscale.

\begin{lstlisting}[style=sqlstyle]
WITH duracion AS (
    SELECT player_id, COUNT(*) AS tics_totales
    FROM TelemetryEvent
    GROUP BY player_id
),
promedio AS (
    SELECT AVG(tics_totales) AS media_global FROM duracion
)
SELECT
    p.player_id,
    p.nickname,
    d.tics_totales,
    s.code AS subescala,
    ROUND(AVG(ri.score), 2) AS puntaje_promedio
FROM duracion d
CROSS JOIN promedio pr
JOIN Player p          ON p.player_id = d.player_id
JOIN UXResponse r      ON r.user_id = p.user_id
JOIN UXResponseItem ri ON ri.response_id = r.response_id
JOIN UXItem i          ON i.item_id = ri.item_id
JOIN UXSubscale s      ON s.subscale_id = i.subscale_id
WHERE d.tics_totales > pr.media_global
GROUP BY p.player_id, p.nickname, d.tics_totales, s.code
ORDER BY p.player_id, s.code;
\end{lstlisting}

\begin{table}[H]
\centering
\begin{tabular}{llr}
\toprule
\textbf{player} & \textbf{nickname} & \textbf{total tics} \\
\midrule
1 & Doomguy & 6000 \\
2 & Ranger & 6300 \\
3 & Crash & 5800 \\
\bottomrule
\end{tabular}
\caption{Players above the global average trajectory duration.}
\end{table}

\textbf{Interpretation.} By combining UX responses with trajectory duration, this query explores whether players who remain active longer report different subjective experiences than average players.

\clearpage
\subsection{Query 5: most visited sector per episode and map}
This query counts visits by episode, map, and sector, then uses \texttt{ROW\_NUMBER()} to keep the most visited sector per map.

\begin{lstlisting}[style=sqlstyle]
WITH conteo AS (
    SELECT
        e.episode_id,
        e.name AS episodio,
        m.map_id,
        m.map_code,
        t.sector_id,
        COUNT(*) AS visitas,
        ROW_NUMBER() OVER (
            PARTITION BY e.episode_id, m.map_id
            ORDER BY COUNT(*) DESC
        ) AS rn
    FROM TelemetryEvent t
    JOIN Game g    ON g.game_id = t.game_id AND g.player_id = t.player_id
    JOIN Map m     ON m.map_id = g.map_id
    JOIN Episode e ON e.episode_id = m.episode_id
    WHERE t.sector_id IS NOT NULL
    GROUP BY e.episode_id, e.name, m.map_id, m.map_code, t.sector_id
)
SELECT episodio, map_code, sector_id AS sector_hotspot, visitas
FROM conteo
WHERE rn = 1
ORDER BY episodio, map_code;
\end{lstlisting}

\begin{table}[H]
\centering
\begin{tabular}{llrr}
\toprule
\textbf{episode} & \textbf{map} & \textbf{hotspot sector} & \textbf{visits} \\
\midrule
Inferno & E3M1 & 246 & 217 \\
Knee-Deep in the Dead & E1M1 & 655 & 293 \\
The Shores of Hell & E2M1 & 489 & 194 \\
\bottomrule
\end{tabular}
\caption{Most visited sector by episode and map.}
\end{table}

\textbf{Interpretation.} Hotspot sectors reveal the areas that attract the highest player activity. These sectors may correspond to combat zones, resource-rich areas, or strategic map locations.

\clearpage
\subsection{Query 6: number of tics where players were together in a sector}
A self-join matches rows with the same \texttt{game\_id}, \texttt{tic}, and \texttt{sector\_id}, but different players. The result counts how many tics each pair shared in the same sector.

\begin{lstlisting}[style=sqlstyle]
SELECT
    a.game_id,
    a.player_id AS jugador_a,
    b.player_id AS jugador_b,
    COUNT(DISTINCT a.tic) AS tics_juntos
FROM TelemetryEvent a
JOIN TelemetryEvent b
  ON a.game_id = b.game_id
 AND a.tic = b.tic
 AND a.sector_id = b.sector_id
 AND a.player_id < b.player_id
GROUP BY a.game_id, a.player_id, b.player_id
ORDER BY tics_juntos DESC;
\end{lstlisting}

\begin{table}[H]
\centering
\begin{tabular}{rrrr}
\toprule
\textbf{game} & \textbf{player A} & \textbf{player B} & \textbf{tics together} \\
\midrule
12 & 4 & 6 & 69 \\
3 & 2 & 6 & 63 \\
4 & 3 & 4 & 60 \\
\bottomrule
\end{tabular}
\caption{Top co-presence results in the same sector.}
\end{table}

\textbf{Interpretation.} Co-presence measurements estimate how often players occupy the same sector simultaneously. High values may indicate collaboration, escort behavior, or repeated encounters.

\section{Indexing and performance evaluation}
\subsection{Important observation before indexing}
The schema already creates implicit B-tree indexes for primary keys and unique constraints. In particular, \texttt{uq\_telem(game\_id, player\_id, tic)} already creates an index over those three columns. Therefore, creating another index with exactly the same columns would be redundant. The evaluated indexes were selected to complement the existing ones.

\subsection{Created indexes}
\textbf{Index 1: \texttt{idx\_telemetry\_game\_tic}.} \\
\textbf{Purpose:} Accelerate proximity and co-presence queries that repeatedly filter telemetry by game and tic.
\begin{lstlisting}[style=sqlstyle]
CREATE INDEX idx_telemetry_game_tic ON TelemetryEvent(game_id, tic);
\end{lstlisting}

\textbf{Index 2: \texttt{idx\_telemetry\_sector}.} \\
\textbf{Purpose:} Improve hotspot and sector-based analyses by avoiding full-table scans when filtering specific sectors.
\begin{lstlisting}[style=sqlstyle]
CREATE INDEX idx_telemetry_sector ON TelemetryEvent(sector_id);
\end{lstlisting}

\textbf{Index 3: \texttt{idx\_game\_player}.} \\
\textbf{Purpose:} Speed up retrieval of all participations associated with a specific player.
\begin{lstlisting}[style=sqlstyle]
CREATE INDEX idx_game_player ON Game(player_id);
\end{lstlisting}

\subsection{Before/after evidence}
\begin{table}[H]
\centering
\footnotesize
\begin{tabular}{p{3.8cm}p{4.2cm}p{2.5cm}p{3.7cm}}
\toprule
\textbf{Index} & \textbf{Representative query} & \textbf{Buffers} & \textbf{Time / plan change} \\
\midrule
\texttt{idx\_telemetry\_sector} & Filter events with \texttt{sector\_id = 655}. & 493 $\rightarrow$ 3 & 2.07 ms $\rightarrow$ 0.069 ms; Seq Scan $\rightarrow$ Index Only Scan. \\
\texttt{idx\_telemetry\_game\_tic} & Proximity inside \texttt{game\_id = 1}. & 168 $\rightarrow$ 150 & 4.9 ms $\rightarrow$ 4.7 ms; Bitmap scan becomes slightly narrower. \\
\texttt{idx\_game\_player} & Search games for \texttt{player\_id = 3}. & 1 $\rightarrow$ 1 & About 0.02 ms $\rightarrow$ 0.02 ms; Seq Scan remains cheaper at 28 rows. \\
\bottomrule
\end{tabular}
\caption{Index evaluation using \texttt{EXPLAIN (ANALYZE, BUFFERS)}.}
\end{table}

\subsection{Discussion}
\textbf{Sector index.} The strongest improvement came from \texttt{idx\_telemetry\_sector}. The filtered sector query changed from a sequential scan over the full telemetry table to an index-only scan. Buffer usage dropped from 493 to 3, and execution time dropped from approximately 2.07 ms to 0.069 ms. This index is especially useful for selective sector queries and co-presence analyses.

\textbf{Game/tic index.} The index \texttt{idx\_telemetry\_game\_tic} had only a marginal improvement in the tested dataset because the unique constraint \texttt{uq\_telem(game\_id, player\_id, tic)} already starts with \texttt{game\_id}. The new index is narrower and slightly reduces buffers, but the difference is small at 34,500 rows.

\textbf{Player index on Game.} The index \texttt{idx\_game\_player} is justified by scalability. With only 28 rows in \texttt{Game}, PostgreSQL correctly chooses a sequential scan because it is cheaper than using the index. However, forcing \texttt{enable\_seqscan = off} confirms that the index is usable, and it will become more valuable as the number of game participations grows.
\clearpage

\section{C.4 Reproducibility mechanism - Makefile}

To ensure full reproducibility of the database, we provide a \texttt{Makefile} that automates the process of resetting and rebuilding the entire system from scratch. This guarantees that the database can be recreated consistently using only the SQL scripts included in the repository.

\subsection{Makefile}

\lstset{
    breaklines=true,
    basicstyle=\ttfamily\small
}

\begin{lstlisting}
DB_NAME=chocolate_doom
DB_USER=postgres

reset:
	psql -U $(DB_USER) -c "DROP DATABASE IF EXISTS $(DB_NAME);"
	psql -U $(DB_USER) -c "CREATE DATABASE $(DB_NAME);"

load: reset
	psql -U $(DB_USER) -d $(DB_NAME) -f "DDL proyecto(1).sql"
	psql -U $(DB_USER) -d $(DB_NAME) -c "\copy TelemetryEvent(game_id, player_id, sector_id, tic, pos_x, pos_y, mom_x, mom_y, angle, fov, health, armor, ammo) FROM 'telemetry.tsv' WITH (FORMAT csv, DELIMITER E'\t', HEADER true);"
	psql -U $(DB_USER) -d $(DB_NAME) -f "C3_Views.sql"
	psql -U $(DB_USER) -d $(DB_NAME) -f "parte2_puntoB.sql"
	psql -U $(DB_USER) -d $(DB_NAME) -f "parte_3_puntoB.sql"
	psql -U $(DB_USER) -d $(DB_NAME) -f "consultas_analiticas.sql"
	psql -U $(DB_USER) -d $(DB_NAME) -f "indices_evaluacion.sql"

views:
	psql -U $(DB_USER) -d $(DB_NAME) -f "C3_Views.sql"

refresh:
	psql -U $(DB_USER) -d $(DB_NAME) -c "REFRESH MATERIALIZED VIEW player_telemetry_summary;"

analytics:
	psql -U $(DB_USER) -d $(DB_NAME) -f "consultas_analiticas.sql"

indexes:
	psql -U $(DB_USER) -d $(DB_NAME) -f "indices_evaluacion.sql"

test:
	psql -U $(DB_USER) -d $(DB_NAME) -c "SELECT * FROM player_game_summary LIMIT 5;"
	psql -U $(DB_USER) -d $(DB_NAME) -c "SELECT * FROM map_activity_summary LIMIT 5;"
\end{lstlisting}

\subsection{Explanation of Makefile targets}

\begin{itemize}

    \item \textbf{reset: }  
    This target completely reinitializes the database environment. It first drops the existing database (if it exists) and then creates a fresh empty instance. This guarantees that every execution starts from a clean and consistent state, avoiding residual data or schema conflicts from previous runs.

    \medskip

    \item \textbf{load: }  
    This is the main orchestration target responsible for fully rebuilding the system from scratch. It first executes \texttt{reset} and then runs all required scripts in a strict dependency order to ensure correctness and referential integrity:
    
    \begin{itemize}
        \item The \textbf{DDL script}, which defines the complete database schema (tables, constraints, and relationships).
        \item The \textbf{TSV ingestion step}, which loads raw telemetry data into the \texttt{TelemetryEvent} table using PostgreSQL's \texttt{\textbackslash copy} mechanism.
        \item The \textbf{C.3 views script}, which creates analytical views and materialized views used for reporting and aggregation.
        \item The \textbf{application logic scripts} (Part 2 and Part 3), which may include additional transformations, constraints, or data processing logic.
        \item The \textbf{analytical queries script}, which defines complex queries used for evaluation and insights extraction.
        \item The \textbf{index evaluation script}, which creates and tests indexes to improve query performance and demonstrates optimization effects.
    \end{itemize}

    This strict ordering ensures that each layer of the system is built on top of a correctly initialized foundation.

    \medskip

    \item \textbf{views: }  
    Recreates only the database views and materialized views defined in the C.3 section without resetting or modifying the underlying data. This is useful when the schema and data already exist but the analytical layer needs to be refreshed or updated.

    \medskip

    \item \textbf{refresh: }  
    Recomputes the materialized view \texttt{player\_telemetry\_summary}. This operation updates precomputed aggregated telemetry results, ensuring that analytical outputs reflect the most recent underlying data.

    \medskip

    \item \textbf{analytics: }  
    Executes the full set of analytical SQL queries defined in \texttt{consultas\_analiticas.sql}. These queries are used to extract insights from the dataset and validate the correctness of the schema and data relationships.

    \medskip

    \item \textbf{indexes: }  
    Executes the index creation and evaluation script defined in \texttt{indices\_evaluacion.sql}. This target is used to measure and demonstrate query performance improvements through indexing strategies.

    \medskip

    \item \textbf{test: }  
    Runs a set of lightweight validation queries over key views. These queries provide a quick sanity check to verify that the database has been correctly built and that the main analytical views are returning expected results.

\end{itemize}
\subsection{Purpose}

This automation allows the database to be fully rebuilt with a single command, without manual steps. It ensures that the setup process is consistent and repeatable, reduces the risk of errors, and makes it easier to reproduce the system in different environments.

\clearpage
\section{Conclusions}
\medskip In this project, we designed and implemented a relational database to store gameplay telemetry together with user information and UX survey data. The schema follows a structured approach that separates core entities such as users, games, telemetry events, and survey responses, making the system organized and scalable.

\medskip One important result is that we can connect player behavior with their subjective experience. By linking telemetry data with survey responses, we can analyze not only how players move in the game, but also how they feel about it. This allows a more complete analysis, for example comparing engagement with movement patterns and activity within the game. Additionally, even though the game is single-player, we extended the analysis to compare spatial behavior across different players.

\medskip We also considered how the data is loaded into the system. By using an ETL approach, we can take raw telemetry logs, clean them, and insert them into the database in a consistent way. This ensures the data is reliable and can be reused if we need to reload everything. The use of staging tables and error logging helps maintain data quality during ingestion.

\medskip Additionally, the use of primary keys, foreign keys, and indexes helps maintain data integrity and improves performance when running analytical queries. This is important because telemetry data grows very fast and queries need to be efficient. Views and materialized views were also implemented to simplify queries and improve performance for frequent analyses.

\medskip Overall, this project shows how a relational database can be used not only to store large amounts of data, but also to support meaningful analysis of gameplay behavior and user experience through real telemetry data and analytical queries.

\appendix
\section{Delivered SQL files}
\begin{itemize}
    \item \texttt{ddl\_corregido.sql}: corrected PostgreSQL schema.
    \item \texttt{parte\_3\_puntoB\_corregido.sql}: UX instrument metadata.
    \item \texttt{populate\_core.sql}: master data and synthetic UX responses.
    \item \texttt{consultas\_analiticas.sql}: six analytical queries.
    \item \texttt{indices\_evaluacion.sql}: index creation and before/after evaluation.
    \item \texttt{evidencia\_explain\_analyze.txt}: raw execution-plan evidence.
\end{itemize}

\end{document}
